% ===========================================================================
% ABSTRACT
% ===========================================================================

\begin{abstract}
Wind turbine speed regulation in above-rated operation requires fast,
constraint-aware control; nonlinear Model Predictive Control (MPC) is
well suited to this, but its internal plant model is expensive to
evaluate at every step, motivating cheaper learned surrogates. We
compare six surrogate architectures --- LLNFM, LSTM, GP, PINN, TCN, and
SW-MLP --- as internal predictors in a nonlinear MPC for the NREL 5-MW
turbine. One-step prediction accuracy and closed-loop tracking prove
only loosely coupled. What initially appeared as a hard operational
incompatibility between LSTM/TCN/SW-MLP/PINN-v2/GP-v2 and SQP-based
\texttt{nlmpc} was traced, through a five-stage diagnostic, to
repeated \texttt{predict()} calls inside the solver --- not
architecture, model size, or MATLAB version. Independently,
\texttt{predict()}'s single-precision output was found to zero out
\texttt{fmincon}'s default finite-difference gradient estimate for
two of these architectures, corrupting solver feasibility regardless
of any real-time deadline. A verified manual forward
pass, computed in double precision, restored closed-loop viability for
all six surrogates. All
results hold under two explicit conditions: a bounded
thirty-iteration SQP solver budget, and a single software toolchain
(MATLAB Deep Learning Toolbox, tested on releases R2024a and R2026a) ---
no cross-toolchain replication was performed. On an Intel Core
i5-12450H (12th Gen, 16\,GB RAM) under R2024a, the headline result is a
$491\times$ mean per-step speedup for LSTM, the largest of the six
architectures tested (full range: $5.5\times$--$491\times$). We further identify a
Region~II/III torque-switching correction as a precondition for stable
closed-loop simulation, and, applying the same diagnostic rigor to our
own results, find that a previously reported "zero pitch activity"
result for a Mat\'ern-kernel GP uncertainty-penalized MPC cost is in
fact a solver-infeasibility artifact (a frozen actuator command, not
genuine regulation) rather than a controller achievement --- a finding
we report as a caution against under-instrumented closed-loop
comparisons, consistent with the diagnostic ethos of this paper.
These results reframe an apparent architectural limitation as a
diagnosable, remediable implementation cost, with practical guidance
for deploying ML surrogates in real-time MPC.

\textbf{Keywords:} Model Predictive Control, Wind Turbine Control,
Machine Learning Surrogates, Gaussian Process, Real-Time Feasibility
\end{abstract}
